\label{chap:conclusion}

This thesis asked whether a U-Net can reproduce algorithmic seabed labels on later EK80 recordings from the same operating area, and whether its bottom estimates can be accompanied by useful uncertainty diagnostics. Within the Bast\o{} Electric route, the answer is cautiously positive. On the held-out Oct 2024 period of 378 files, none seen in training, the Bast\o{} Standard U-Net reached a mean IoU of 0.9972 and 0.57\,m mean depth deviation against the reference detector. The same checkpoint retained a mean IoU of 0.998 and 0.516\,m mean depth deviation on the later 135-file Feb--Jun 2025 test split, with finite depth deviation for 134 files. These results establish strong within-route model-reference agreement across recording periods, not geographic transfer or physical seabed accuracy.

The uncertainty results are the main practical gain over a traditional detector that returns only a point estimate. The raw sigmoid output is useful as a boundary diagnostic, because ambiguous pixels concentrate near the seabed transition, but it is not itself a reliable uncertainty guarantee: the probability field is nearly binary, the expected calibration error is low partly for that reason, and the reliability diagram still shows miscalibration. Split-conformal calibration gives the more reportable supervised interval, but its success criterion is reference-label coverage, not true-seabed coverage. At the 90\% target, the supervised constant-width band calibrated on the January 2025 cache achieved 92.6\% pooled empirical reference-label coverage on the scored Feb and Jun 2025 conformal cache. This pooled result does not establish general within-route calibration: coverage was 95.6\% in February and 89.8\% in June, varied materially across files, and the calibration and test score distributions differed. A second, label-free construction from repeated crossings evaluates pairwise route repeatability rather than physical seabed coverage or per-ping accuracy. Both bands cover held-out crossing pairs at about 90.2\%, while the adaptive band cuts the median half-width from 7.59\,m to 3.59\,m by widening mainly where crossings disagree. Patch-scale examples show the same behaviour on individual echograms. The band contracts on easy flat bottoms, enough to miss an essentially correct prediction, and widens where the model and reference diverge, so its width flags locally uncertain seabed estimates, by eye or as an automatic threshold, rather than providing a calibrated per-ping interval.

The dominant weakness is the reference. All supervised metrics measure agreement with algorithmic labels, not distance from the true seabed, and this includes supervised conformal coverage when the nonconformity score is computed against those labels. The labels are generated from 38\,kHz detections, and the channel ablation shows that the trained network also depends mainly on 38\,kHz, so a high score can mean that the model has learned the detector's response rather than an independent seabed representation. Selective visual inspection suggests that some large model-reference errors come from detector and labelling faults, but this remains visual evidence, not an independently labelled error rate. Repeated crossings expose the same problem without using labels: at fixed route bins the assigned depth varies by tens of metres across passes, especially in the deep central channel and adjacent slope. The harbour approaches add a separate failure mode, where shallow water, braking, and bow-thruster bubbles make the labels unreliable in a zone where operational errors matter.

The second major weakness is limited geographic and operational diversity. Nearly all training, calibration, and within-route testing comes from one ferry, one transducer installation, and one repeated bottom profile. The chronological split tests later acoustic conditions on the same route, but it does not test unseen seabed, other vessels, different pulse settings, or different biological regimes. The NOAA Pacific hake analyses are therefore best read as limited stress tests rather than broad external validation. In the zero-shot diagnostic, the unchanged Bast\o{} Standard U-Net is evaluated on six selected stitched hake echograms, where agreement ranges from 0.61 to 0.84 IoU and dense pelagic scattering layers are sometimes marked as bottom. This range describes those six examples, not the full external dataset. A separate threshold-only recalibration uses 20 hake files for calibration and 49 held-out files for evaluation. It restores constant-width reference coverage to 89.5\%, near the 90\% target, but expands the median band width to 42.3\,m and still measures coverage against algorithmic labels. Model adaptation uses a third file-level split with 60 training files, 15 validation files, 9 calibration files, and 9 held-out test files. On the three stitched groups formed from the nine adaptation test files, fine-tuning raises mean IoU from 0.433 to 0.683, reduces mean depth deviation from 30.11 to 12.69\,m, and increases Bast\o{}-calibrated constant-width coverage from 7.9\% to 89.8\%. A Hake-only scratch model performs better on these small adaptation test groups but collapses on Bast\o{}. The three hake protocols are separate diagnostic runs over the processed pool, and their file counts are not additive. The results support domain-specific deployment and separate validation of adapted checkpoints, but they do not establish that successful transfer is possible only through specialisation or that an adapted model will generalise to other surveys. All supervised results remain relative to algorithmic reference labels.

Encoder strategy, loss choice, and dataset variants were compared on the 89-file Jan 2024 development subset, using an 80/20 file-level split within the same recording period. Within the encoder family, the freezing schedule made only small differences, the larger ResNet-50 backbone did not improve on ResNet-34, and ImageNet pretraining was not required for high within-route agreement. They were not evaluated after full-data training, and their small IoU differences should be interpreted as controlled development results rather than temporal generalisation effects. A separate five-seed deep ensemble, evaluated on the full Oct 2024 and Feb--Jun 2025 data rather than the development subset, did not improve point accuracy over a single member. The augmentation ablation used a separate 100-file random subset, with 80 files for training and 20 for validation. Within that experiment, channel dropout and ping-stripe corruption reduced depth deviation and across-file spread, consistent with encouraging the model to use spatial context rather than one intact channel. Since the two ablation families use different subsets, the results do not establish that augmentation is more influential than encoder strategy, loss, or dataset construction. The channel study points to a simpler future model, since 38\,kHz carries most of the seabed-placement information and the present difference channel contributes little under the current normalisation.

The next step is therefore not mainly a larger network, but a better reference, broader data, and deployment-aware training. Expert-labelled subsets, detector-disagreement filters, and repeated-crossing quality metrics would separate locations where the reference is trustworthy from locations where training should be down-weighted, corrected, or excluded. A live system should also match the causal deployment geometry: when the newest ping sits at the right edge of a patch with no future context, the final column degrades, so the model should either train in that geometry, report with a short delay, or validate a mirroring strategy directly. Fully empty recordings also need a separate absent-bottom evaluation rather than inclusion in the bottom-bearing headline metrics. The thesis contributes a bounded operational case study of learned EK80 seabed segmentation on a fixed route. It shows that a U-Net can reproduce a reference detector and provide empirical diagnostics of reference-label discrepancy and route repeatability. The sigmoid output was not operationally calibrated, and dependence and distribution shift limit the validity of the supervised conformal results. Broader use therefore requires independent seabed references, more geographically diverse validation, and uncertainty checks under distribution shift.
